\scanpage{90}
\mspage{9}

action tautologique de $\unsure{\mathfrak S|\mathcal G}_3$, et sur $\mathord{\mathrm{I}\mkern-7mu\mu}_4^*$ via le caractère mixte $\operatorname{sg}(g)\cdot\deleted{\unsure{\zeta}}\lambda^2$ $(g\in \mathfrak S_3,\deleted{\zeta}\lambda\in \mathord{\mathrm{I}\mkern-7mu\mu}_4)$. Donc le sous-groupe de $\mathfrak S_3\times \mathord{\mathrm{I}\mkern-7mu\mu}_4$ qui opère trivialement sur $\mathord{\mathrm{I}\mkern-7mu\mu}_4^*$ est

\[
\mathbb D'_{3S_0}=(\mathfrak S_3)_{S_0}\times_{\deleted{S}\mathord{\mathrm{I}\mkern-7mu\mu}_{2S_0}}(\mathord{\mathrm{I}\mkern-7mu\mu}_4)_{S_0}
\tag*{(48)}
\]

Si on prend comme base absolue non $S_0$, mais $\widetilde S_0=(\mathord{\mathrm{I}\mkern-7mu\mu}_4^*)_{S_0}$ de sorte que l'on trouve, sur les schémas \inserted{$/\widetilde S_0$} envisagés

\[
\left\{\begin{aligned}
&\mathord{\mathrm{I}\mkern-7mu\mu}_{4\deleted{\unsure{0}}S}^*\ \text{muni d'une section canonique }i_S\\
&\mathord{\mathrm{I}\mkern-7mu\mu}_{4S}\simeq(\mathbb Z/4\mathbb Z)_S
\end{aligned}\right.
\]

alors le groupe $(48)$ s'identifie au groupe constant sur $S$ \deleted{de valeur} \inserted{défini par le groupe \unsure{concret|ensembliste|constant|ordinaire}}

\[
\mathbb D'_3=\mathfrak S_3\times_{\mathbb Z/2\mathbb Z}(\mathbb Z/4\mathbb Z)
\tag*{(49)}
\]

qui est une extension de $\mathbb D_3\simeq \mathfrak S_3$ par $\mathbb Z/2\mathbb Z\simeq\{\pm1\}$

\[
1\mathbin{\relbar}\deleted{\mathbb Z/2}\{\pm1\}\mathbin{\relbar}\mathbb D'_3\mathbin{\relbar}\mathbb D_3\mathbin{\relbar}1\quad .
\tag*{(50)}
\]

\deleted{\unsure{Jusqu'ici|Change ici}} \deleted{L'action envisagée de $\mathfrak S_3\times \mathord{\mathrm{I}\mkern-7mu\mu}_4$} \inserted{Dans le contexte i.e. en} travaillant sur $\widetilde S$ i.e. avec un choix fixé de $i=\sqrt{-1}$, on peut distinguer

% archive sheet number: 89
